%(BEGIN_QUESTION)
% Copyright 2006, Tony R. Kuphaldt, released under the Creative Commons Attribution License (v 1.0)
% This means you may do almost anything with this work of mine, so long as you give me proper credit

Jeg kan fylle et kar på 4 liter fra kjøkkenkranen på 30 sekunder. Røret til kjøkkenkranen har ID = 12.5 mm. Beregn Reynolds-tallet for denne strømningen. Blir strømningen turbulent eller laminær?

\underbar{file i00442}
%(END_QUESTION)





%(BEGIN_ANSWER)

Re $\approx$ 13581 = {\it turbulent}

\vskip 10pt

Reynolds-tall under 2,000 tilsvarer vanligvis laminær strømning, mens Reynolds-tall over 10,000 vanligvis tilsvarer turbulent strømning. Mellom 2,000 og 10,000 er strømningen ofte svakt turbulent, gjerne kalt «overgangsstrømning». Merk at disse grensene er {\it omtrentlige} og avhenger av blant annet rørgeometri og ruhet i veggen.

Eksempler på publiserte terskelverdier for laminær/turbulent strømning:

\begin{itemize}
\goodbreak
\item{} Re $<$ 2,000 = ``Laminær''
\item{} 2,000 $<$ Re $<$ 10,000 = ``Overgang''
\item{} Re $>$ 10,000 = ``Fullt utviklet turbulent''
\item{} Kilde: R. Siev, J.B. Arant, B.G. Lipt\'ak; \underbar{Chapter 2.8: Laminar Flowmeters}; {\it Instrument Engineer's Handbook, Process Measurement and Analysis, Third Edition}; pg. 105
\end{itemize}

\begin{itemize}
\goodbreak
\item{} Re $>$ 10,000 = ``Definitivt turbulent''
\item{} Kilde: W.H. Howe, J.B. Arant, B.G. Lipt\'ak; \underbar{Chapter 2.14: Orifices}; {\it Instrument Engineer's Handbook, Process Measurement and Analysis, Third Edition}; pg. 153
\end{itemize}

\begin{itemize}
\goodbreak
\item{} Re $<$ 2,000 = ``Laminær''
\item{} 2,000 $<$ Re $<$ 4,000 = ``Overgang''
\item{} Re $>$ 4,000 = ``Turbulent''
\item{} Kilde: Instrument Society of America; \underbar{Chapter 2: Fluid Properties -- Part II}; {\it ISA Industrial Measurement Series -- Flow}; pg. 11
\end{itemize}

\begin{itemize}
\goodbreak
\item{} Re $<$ 2,100 = ``Laminær''
\item{} Re $>$ 3,000 = ``Turbulent''
\item{} Kilde: Tyler G. Hicks, P.E.; \underbar{Laminar Flow in a Pipe}; {\it Standard Handbook of Engineering Calculations}; pg. 1-202
\end{itemize}

\begin{itemize}
\goodbreak
\item{} Re $<$ 1,200 = ``Laminær''
\item{} Re $>$ 2,500 = ``Turbulent''
\item{} Kilde: Tyler G. Hicks, P.E.; \underbar{Piping and Fluid Flow}; {\it Standard Handbook of Engineering Calculations}; pg. 3-384
\end{itemize}

Det er litt komisk hvor forskjellige terskelverdier som oppgis i litteraturen.

\begin{itemize}
\goodbreak
\item{} Re $<$ (omtrent) 2,000 = ``Laminær''
\item{} Re $>$ 2,000 = ``Turbulent''
\item{} Kilde: Douglas C. Giancoli; \underbar{Chapter 10: Fluids}; {\it Physics (Third Edition)}; pg. 11
\end{itemize}

\begin{itemize}
\goodbreak
\item{} Re $<$ (omtrent) 2,000 = ``Laminær''
\item{} 2,000 $<$ Re $<$ 4,000 = ``Overgang''
\item{} Re $>$ 4,000 = ``Turbulent''
\item{} Kilde: Schoolcraft Publishing; \underbar{Chapter 20: Properties of Fluid Flow}; {\it Process Instrumentation -- Volume I}; pg. 258
\end{itemize}

\goodbreak

En annen kilde, nesten komisk presis i avgrensningen, oppgir:

\begin{itemize}
\item{} Re $<$ 2,320 = ``Laminær''
\item{} Re $>$ 2,320 = ``Turbulent''
\item{} Kilde: Nettside ({\tt http://flow.netfirms.com/reynolds/theory.htm})
\end{itemize}

Det bør også nevnes at laminær strømning kan opprettholdes ved Reynolds-tall klart over 10,000 under spesielle forhold. For eksempel i enkelte kveilede kapillarrør kan laminær strømning opprettholdes helt opp mot Re = 15,000, på grunn av det som kalles {\it Dean-effekten}.

%(END_ANSWER)





%(BEGIN_NOTES)

$$Q=\frac{4l}{30s}=0.13 l/s$$
$$V=\frac{0.00013m^3/s}{\pi \frac{0.0125m^2}{4}}=1.1m/s$$
$$Re=\frac{\rho vD}{\mu}=\frac{1000kg/m^3 \cdot 1.1m/s \cdot 0.0125m}{1m Pa\cdot s}=13581 \hbox{ (turbulent)}$$
%INDEX% Physics, dynamic fluids: Reynolds number

%(END_NOTES)
